\documentclass[11pt,a4paper,twocolumn]{article}

% === Core Packages ===
\usepackage[margin=2cm]{geometry}
\usepackage[utf8]{inputenc}
\usepackage[T1]{fontenc}
\usepackage{times}
\usepackage{amsmath,amssymb}
\usepackage{graphicx}
\usepackage{booktabs}
\usepackage{float}
\usepackage{enumitem}
\usepackage{tikz}
\usepackage{pgfplots}
\pgfplotsset{compat=1.18}
\usetikzlibrary{arrows.meta,shapes.geometric,positioning,calc,decorations.markings,fit,backgrounds}

% === Citation (IEEE style) ===
\usepackage[numbers,sort&compress]{natbib}
\bibliographystyle{IEEEtranN}

% === Hyperref ===
\usepackage[hidelinks]{hyperref}
\usepackage{xurl}

% === Settings ===
\setlength{\parindent}{1em}
\setlength{\parskip}{0.3em}

% === PDF metadata ===
\hypersetup{
  pdftitle={เครือข่ายชื่อเสียงแบบกระจายศูนย์: กรอบแนวคิดสำหรับการจัดระเบียบสังคมตามธรรมชาติ},
  pdfauthor={Pavel Kudrna},
  pdfsubject={Personix — a decentralized reputation network},
  pdfkeywords={Personix, decentralized reputation, reputation social network,
    decentralized identifiers, DID, meritocracy, censorship resistance},
}

% === Title ===
\title{\textbf{เครือข่ายชื่อเสียงแบบกระจายศูนย์:\\กรอบแนวคิดสำหรับการจัดระเบียบสังคมตามธรรมชาติ}}
\author{Pavel Kudrna\\Personix Project --- \url{https://personix.org}\\Prague, Czechia\\ฉบับร่าง v1 --- มีนาคม 2026}
\date{}

\begin{document}
\maketitle

% ============================================================
% ABSTRACT
% ============================================================
\begin{abstract}
สำนึกแห่งความยุติธรรม---ความเชื่อมั่นว่าความผิดจะได้รับการแก้ไขให้ถูกต้องและคุณความดีจะได้รับการตอบแทน---คือแรงยึดเหนี่ยวที่ทรงพลังที่สุดของสังคม เมื่อสถาบันการปกครองแบบรวมศูนย์ไม่สามารถมอบสำนึกนี้ได้ เพราะต้นทุนของการบังคับใช้สิทธิหนึ่งสูงเกินกว่ามูลค่าของสิทธินั้นเอง ความเป็นน้ำหนึ่งใจเดียวกันของสังคมก็เสื่อมสลาย โครงสร้างของสถาบันในปัจจุบันล้มเหลวในการระงับข้อพิพาทจำนวนมากอย่างเป็นระบบ ในลักษณะเดียวกับที่การวางแผนจากส่วนกลางล้มเหลวในการจัดสรรทรัพยากร นั่นคือ ผ่านความไม่สมมาตรของข้อมูล การขาดวงจรป้อนกลับ และการที่ผู้ตัดสินใจถูกตัดขาดจากผลลัพธ์ของการตัดสินใจของตน เอกสารฉบับนี้นำเสนอเครือข่ายสังคมชื่อเสียง (Reputation Social Network, RSN) ซึ่งเป็นชั้นการสื่อสารแบบกระจายศูนย์ ต้านทานการทุจริต และต้านทานการเซ็นเซอร์ สร้างขึ้นบนตัวระบุแบบกระจายศูนย์ (Decentralized Identifiers, DIDs) ที่เปิดโอกาสให้ผู้เข้าร่วมแบบนามแฝงสามารถเผยแพร่ ตรวจสอบ และบริโภคข้อมูลชื่อเสียงเกี่ยวกับพฤติกรรมในโลกจริงได้ กลไกหลักตั้งอยู่บนหลักการออกแบบสามประการ ได้แก่ (1)~โครงสร้างต้นทุนแบบไม่สมมาตร ที่การอ่านข้อมูลชื่อเสียงมีราคาถูก แต่การเขียนต้องอาศัยความพยายามที่พิสูจน์ได้; (2)~โปรโตคอลการคัดเลือกผู้ตรวจสอบแบบไม่กำหนดตายตัว ที่อาศัยการแฮชแบบสม่ำเสมอ (consistent hashing) ซึ่งตั้งราคาข้ออ้างสุดโต่งผ่านต้นทุนการทำซ้ำที่เพิ่มสูงขึ้นเรื่อยๆ; และ (3)~แบบจำลองผู้มีอำนาจด้านชื่อเสียงที่ขับเคลื่อนด้วยตลาด ซึ่งผู้ตรวจสอบเอกชนเดิมพันชื่อเสียงที่สั่งสมมาไว้กับความถูกต้องของข้ออ้างที่ตนรับรอง สถาปัตยกรรมที่เกิดขึ้นก่อให้เกิดระบบคุณธรรมนิยม (meritocracy) แบบผุดขึ้นเองโดยไม่ต้องมีการประสานงานจากส่วนกลาง และเปิดเส้นทางการโยกย้ายที่ค่อยเป็นค่อยไปและย้อนกลับได้จากระบบที่รัฐบริหารอยู่ในปัจจุบัน
\end{abstract}

% ============================================================
% 1. INTRODUCTION
% ============================================================
\section{บทนำ}

\subsection{ปัญหาของความไว้วางใจแบบรวมศูนย์}

การปกครองสมัยใหม่ตั้งอยู่บนสมมติฐานพื้นฐานประการหนึ่ง นั่นคือ สถาบันแบบรวมศูนย์สามารถเป็นตัวกลางแห่งความไว้วางใจระหว่างคนแปลกหน้าในระดับกว้างขวางได้ ศาลตัดสินข้อพิพาท ทะเบียนรับรองกรรมสิทธิ์ องค์กรกำกับดูแลบังคับใช้มาตรฐาน สถาบันเหล่านี้ถูกออกแบบขึ้นในยุคที่ข้อมูลหายาก การสื่อสารเป็นไปอย่างเชื่องช้า และการมอบอำนาจให้แก่หน่วยงานส่วนกลางเป็นกลไกการประสานงานเพียงอย่างเดียวที่เป็นไปได้ อย่างไรก็ตาม การปกครองแบบรวมศูนย์ล้มเหลวด้วยเหตุผลเชิงโครงสร้างเดียวกันกับการวางแผนจากส่วนกลาง นั่นคือ ความไม่สมมาตรของข้อมูล การขาดวงจรป้อนกลับ และการที่ผู้ตัดสินใจถูกตัดขาดจากผลลัพธ์ของการตัดสินใจของตน

สมมติฐานนี้ล้มเหลว ยกตัวอย่างเช่น เมื่อต้นทุนของการไกล่เกลี่ยโดยสถาบันสูงเกินกว่ามูลค่าของข้อพิพาท ลองพิจารณาผู้เช่ารายหนึ่งที่พบว่าอพาร์ตเมนต์ที่ตนเช่าอยู่ขาดใบรับรองความปลอดภัยทางไฟฟ้าตามที่กฎหมายกำหนด---สภาพที่ก่อให้เกิดอันตรายต่อชีวิตในทันที สิทธิตามกฎหมายของผู้เช่าในการถอนตัวจากสัญญานั้นชัดเจนไม่คลุมเครือ แต่ต้นทุนของการบังคับใช้สิทธินั้นผ่านระบบศาลกลับสูงเกินกว่ามูลค่าเป็นตัวเงินของเงินมัดจำและค่าเสียหายที่ต้องชดใช้ แม้จะรวมค่าชดเชยตามกฎหมายที่อาจได้รับไว้ด้วยแล้วก็ตาม~\cite{czcourtfees} การตอบสนองอย่างมีเหตุผลคือการยอมรับความสูญเสียและถอนตัวออกไป

นี่ไม่ใช่กรณีเฉพาะที่ผิดปกติ แต่คือประสบการณ์ตามปกติสำหรับข้อพิพาทจำนวนมาก ที่ความเสียหายเกิดขึ้นจริง แต่มูลค่าเป็นตัวเงินต่ำกว่าเกณฑ์ที่การบังคับใช้โดยสถาบันจะคุ้มค่าทางเศรษฐกิจ ผลลัพธ์คือระบบที่เอื้อประโยชน์เชิงโครงสร้างแก่ผู้กระทำผิด นั่นคือ ผู้ที่ทำร้ายผู้อื่นในปริมาณต่ำกว่าเกณฑ์นี้สามารถกระทำได้โดยไม่ต้องรับผิด เพราะต้นทุนของการแสวงหาความยุติธรรมตกอยู่กับฝ่ายผู้เสียหาย

ตัวอย่างข้างต้นดึงมาจากสภาพแวดล้อมทางกฎหมายของผู้เขียน---ระบบกฎหมายแพ่งของเช็ก ในประเพณีทางกฎหมายอื่นๆ---กฎหมายจารีตประเพณีที่อิงบรรทัดฐาน (common law) กฎหมายจารีตประเพณี ระบบผสม---ความยุติธรรมล้มเหลวในลักษณะที่แตกต่างกันและในระดับที่แตกต่างกัน ขึ้นอยู่กับลักษณะเฉพาะทางวัฒนธรรมและกระบวนพิจารณาของเขตอำนาจศาลนั้นๆ ผู้อ่านมีแนวโน้มที่จะรับรู้ถึงตัวอย่างที่เทียบเคียงกันได้จากบริบทของตนเอง สิ่งที่เป็นสากลในเชิงโครงสร้างคือรูปแบบนี้ นั่นคือ ข้อพิพาทที่มูลค่าต่ำกว่าเกณฑ์ของการบังคับใช้ที่คุ้มค่าทางเศรษฐกิจจะจบลงด้วยการที่ฝ่ายผู้เสียหายต้องแบกรับความสูญเสียเอง RSN มิได้สัญญาว่าจะมอบความยุติธรรมอันสมบูรณ์แบบ---เพียงแต่ลดความได้เปรียบเชิงโครงสร้างที่ผู้กระทำผิดได้รับ เมื่อการบังคับใช้โดยสถาบันไม่คุ้มค่าทางเศรษฐกิจ

\subsection{เหตุใดวิธีแก้ปัญหาที่มีอยู่จึงยังไม่เพียงพอ}

\textbf{กรอบแนวคิดอัตลักษณ์แบบกระจายศูนย์ (DID)}~\cite{did2022} มอบกลไกทางวิทยาการเข้ารหัสสำหรับการจัดการอัตลักษณ์แบบอธิปไตยในตนเอง (self-sovereign identity) แต่มุ่งเน้นหลักไปที่การตรวจสอบอัตลักษณ์ โดยไม่ได้กล่าวถึงคำถามที่กว้างกว่าเกี่ยวกับชื่อเสียงด้านพฤติกรรม

\textbf{โทเคนผูกวิญญาณ (Soulbound Tokens, SBTs)}~\cite{weyl2022} จับใจความสำคัญที่ว่าชื่อเสียงเป็นสิ่งที่โอนกันไม่ได้ แต่กลับพึ่งพาโครงสร้างพื้นฐานบล็อกเชนซึ่งมีข้อจำกัดด้านการขยายขนาด และไม่ได้กล่าวถึงแรงจูงใจทางเศรษฐกิจที่กำกับการป้อนข้อมูลเข้าสู่ระบบ แนวคิดที่เกี่ยวข้อง เช่น โทเคนคุณความดี (merit tokens) (เช่นของ Liberland) มีปรัชญาที่คล้ายคลึงกัน แต่ยังคงผูกติดอยู่กับกรอบเขตอำนาจศาลเฉพาะแห่ง

\textbf{องค์กรอิสระแบบกระจายศูนย์ (DAOs)}~\cite{buterin2014} ดำเนินการปกครองผ่านสัญญาอัจฉริยะ (smart contracts) แต่กลับจำลองพลวัตแบบธนาธิปไตย (plutocratic) ที่อิทธิพลเป็นสัดส่วนกับทุน การลงคะแนนแบบกำลังสอง (quadratic voting)~\cite{buterin2019} ช่วยบรรเทาปัญหานี้ได้บางส่วน แต่ก็จำกัดการนำไปใช้ในทางปฏิบัติ นอกจากนี้ DAOs ยังเผชิญกับ\emph{ปัญหาออราเคิล} (oracle problem) อันเป็นพื้นฐาน นั่นคือ สัญญาอัจฉริยะไม่สามารถตรวจสอบสถานะในโลกจริงได้อย่างน่าเชื่อถือหากปราศจากแหล่งข้อมูลภายนอกที่ไว้ใจได้ และปัญหานี้ก็ไม่มีทางแก้ทั่วไปภายในสถาปัตยกรรมบล็อกเชน

ไม่มีระบบใดที่มีอยู่ในปัจจุบันที่ผสานการกระจายศูนย์ ความต้านทานการเซ็นเซอร์ ความเป็นนามแฝง ต้นทุนการเข้าสู่ระบบที่พิสูจน์ได้ และฉันทามติแบบผุดขึ้นเองไว้ด้วยกัน

\subsection{คุณูปการ}

เอกสารฉบับนี้มีคุณูปการดังต่อไปนี้:
\begin{enumerate}[nosep]
\item การนิยาม\textbf{เครือข่ายสังคมชื่อเสียง (RSN)} ซึ่งเป็นโปรโตคอลแบบกระจายศูนย์ที่ขยายกรอบแนวคิด DID เพื่อรองรับการแลกเปลี่ยนชื่อเสียงแบบสองทาง
\item \textbf{โปรโตคอลการคัดเลือกผู้ตรวจสอบแบบไม่กำหนดตายตัว} ที่อาศัยการแฮชแบบสม่ำเสมอ~\cite{karger1997}
\item \textbf{หลักการความสุดโต่งราคาแพง (Expensive Radicalism Principle)} ที่ตั้งราคาข้ออ้างตามระยะห่างจากฉันทามติของเครือข่าย ผ่านช่องทางต้นทุนที่ทบต้นกันสามช่องทาง
\item \textbf{แบบจำลองผู้มีอำนาจที่น่าเชื่อถือ (Reputable Authority model)} ที่มีแรงจูงใจแบบไม่สมมาตร ซึ่งการรับรองอันเป็นเท็จมีต้นทุนสูงกว่าค่าธรรมเนียมที่ชอบด้วยกฎหมายอย่างมหาศาล
\item \textbf{เส้นทางการโยกย้ายที่ค่อยเป็นค่อยไป} ผ่านโครงสร้างพื้นฐานทางการคลังแบบคู่ขนาน
\end{enumerate}

% ============================================================
% 2. DESIGN PRINCIPLES
% ============================================================
\section{หลักการออกแบบ}

สถาปัตยกรรมของ RSN ถูกกำกับด้วยหลักการออกแบบที่ต่อรองไม่ได้ห้าประการ

\textbf{ความต่อเนื่องและวิวัฒนาการ} ระบบอยู่ร่วมกับสถาบันที่มีอยู่เดิมในช่วงเปลี่ยนผ่านที่มีระยะเวลาไม่แน่นอน การเปลี่ยนผ่านต้องย้อนกลับได้ในทุกขั้นตอน

\textbf{ความสมัครใจและความรับผิดชอบ} การเข้าร่วมเป็นไปโดยสมัครใจ แต่ความสมัครใจถูกเชื่อมโยงกับความรับผิดชอบ นั่นคือ ผู้เข้าร่วมที่รับหน้าที่ควบคุมบทบาทของสถาบันหนึ่งย่อมรับภาระผูกพันที่มาพร้อมกันไปด้วยในเวลาเดียวกัน

\textbf{ความหลากหลายและความเป็นกลางทางวัฒนธรรม} ฉันทามติผุดขึ้นจากปฏิสัมพันธ์แบบสองทาง มิใช่จากกฎเกณฑ์ที่ตั้งไว้ล่วงหน้าซึ่งกำหนดโดยผู้ออกแบบระบบ

\textbf{ความยืดหยุ่นและความต้านทานการเซ็นเซอร์} ต้นทุนของการเซ็นเซอร์เครือข่ายต้องสูงกว่าต้นทุนของการยอมทนต่อมัน มีเพียงระบอบเผด็จการเบ็ดเสร็จอย่างสมบูรณ์เท่านั้น---ด้วยต้นทุนทางการเมืองและเศรษฐกิจอันมหันต์---ที่จะสามารถกดขี่ระบบนี้ได้

\textbf{ฉันทามติและการบังคับใช้} ระบบผนวกแนวโน้มของมนุษย์ในด้านความใจแคบ ความอาฆาต และความสะใจในการแก้แค้นเข้าไว้เป็นคุณลักษณะรับน้ำหนัก การประพฤติผิดไม่ได้ถูกห้าม แต่ถูกตั้งราคา

% ============================================================
% 3. SYSTEM OVERVIEW
% ============================================================
\section{ภาพรวมของระบบ}

\subsection{เครือข่ายสังคมชื่อเสียง}

RSN คือชั้นการสื่อสารแบบเพียร์ทูเพียร์ (peer-to-peer) ที่กระจายศูนย์ ซึ่งผู้เข้าร่วมแลกเปลี่ยนข้อมูลที่ตรวจสอบได้เกี่ยวกับพฤติกรรมในโลกจริง คุณสมบัติที่นิยามตัวมันได้แก่ โครงสร้างพื้นฐานที่กระจายศูนย์และเซ็นเซอร์ไม่ได้; การเข้าร่วมแบบนามแฝงผ่าน DID; โครงสร้างต้นทุนแบบไม่สมมาตร (การอ่านมีราคาถูก การเขียนมีราคาแพง); ความสามารถในการตรวจสอบทางวิทยาการเข้ารหัส; และ\textbf{การรองรับมาตรฐาน}---รูปแบบที่เครื่องอ่านได้สำหรับข้อมูลที่แพร่กระจายผ่านเครือข่าย สำหรับบริการที่ผู้มีอำนาจนำเสนอ (การประกอบข้ออ้าง การรับรอง บริการเป็นพยาน) และสำหรับรูปแบบนโยบายซึ่งรวมถึงกฎเกณฑ์สำหรับการประเมินชื่อเสียงของคู่สัญญา และการประเมินความเสี่ยงของความไม่สอดคล้องของนโยบาย (ระหว่างพฤติกรรมที่ประกาศไว้กับพฤติกรรมที่แท้จริง)

\subsection{องค์ประกอบทางสถาปัตยกรรม}

RSN ประกอบด้วยองค์ประกอบหลักสี่ส่วน (รูปที่~\ref{fig:architecture}):
\begin{enumerate}[nosep]
\item \textbf{ชั้น DID} อัตลักษณ์ของผู้เข้าร่วมพร้อมกฎที่ประกาศไว้ ข้อมูลอภิพันธุ์ชื่อเสียง และการกำหนดเส้นทางในเครือข่าย
\item \textbf{โปรโตคอลข้ออ้าง} รูปแบบที่มีโครงสร้างสำหรับการเผยแพร่คำยืนยันเกี่ยวกับเหตุการณ์ในโลกจริง
\item \textbf{เครือข่ายการตรวจสอบ} โปรโตคอลแบบกระจายศูนย์สำหรับการมอบหมายผู้ตรวจสอบโดยใช้การแฮชแบบสม่ำเสมอ
\item \textbf{ชั้นการรวบรวมชื่อเสียง} บริการที่ผลิตบทสรุปชื่อเสียงที่มนุษย์อ่านได้
\end{enumerate}

\begin{figure}[ht]
\centering
\begin{tikzpicture}[
    layer/.style={draw, minimum height=0.7cm, align=center, font=\scriptsize, fill=black!8},
    arr/.style={<->, thick, gray!60}
]
\node[layer, minimum width=5cm] (agg) at (0,2.8) {\textbf{การรวบรวม}\\ตีความ $\cdot$ ความเสี่ยง};
\node[layer, minimum width=2.2cm] (claim) at (-1.55,1.4) {\textbf{ข้ออ้าง}\\ประกอบ};
\node[layer, minimum width=2.2cm] (verif) at (1.55,1.4) {\textbf{การตรวจสอบ}\\การคัดเลือก};
\node[layer, minimum width=5cm] (did) at (0,0) {\textbf{ชั้น DID}\\อัตลักษณ์ $\cdot$ กฎ $\cdot$ คะแนน};

\draw[arr] (agg.south west) -- (claim.north);
\draw[arr] (agg.south east) -- (verif.north);
\draw[arr] (claim.east) -- (verif.west);
\draw[arr] (claim.south) -- (did.north west);
\draw[arr] (verif.south) -- (did.north east);
\end{tikzpicture}
\caption{สถาปัตยกรรมสี่ชั้นของ RSN}
\label{fig:architecture}
\end{figure}

\subsection{บทบาทของผู้เข้าร่วม}

ธุรกรรมการตรวจสอบหนึ่งเกี่ยวข้องกับบทบาท DID ที่แตกต่างกันได้มากถึงหกบทบาท (รูปที่~\ref{fig:roles}): \textbf{ผู้ออกข้ออ้าง} ($DID_I$ สร้างและส่งข้ออ้าง), \textbf{ผู้เป็นเป้าหมายข้ออ้าง} ($DID_S$ คือ DID ที่ข้ออ้างนั้นกล่าวถึง---อาจเป็นคนเดียวกับผู้ออกข้ออ้าง), \textbf{ผู้มีอำนาจ} ($DID_A$ รับรองข้ออ้างโดยมีค่าธรรมเนียม; อาจเสนอบริการเป็นพยานด้วย---\emph{เป็นทางเลือก}), \textbf{พยาน} ($DID_{W_{1..n}}$ ผู้เฝ้าสังเกตความซื่อสัตย์ของผู้ตรวจสอบแบบนิรนามผ่านรหัสท้าทาย---\emph{เป็นทางเลือก}), \textbf{ผู้ตรวจสอบ} ($DID_V$ ถูกคัดเลือกโดยอัลกอริทึมเพื่อเผยแพร่ข้ออ้าง) และ \textbf{RSN} (เครือข่ายที่ข้ออ้างซึ่งผ่านการตรวจสอบแล้วถูกเผยแพร่) บทบาทใดๆ อาจถูกมอบหมายให้ DID อื่นได้ (หัวข้อ 7.4) ผู้มีอำนาจมักรวมการรับรองเข้ากับบริการเป็นพยาน แต่ทั้งสองหน้าที่นี้เป็นอิสระต่อกันในเชิงตรรกะ

\begin{figure}[ht]
\centering
\begin{tikzpicture}[
    role/.style={draw, rounded corners, minimum width=1.1cm, minimum height=0.45cm, align=center, font=\scriptsize, fill=black!8},
    opt/.style={draw, rounded corners, minimum width=1.1cm, minimum height=0.45cm, align=center, font=\scriptsize, fill=black!8, dashed},
    arr/.style={-{Stealth[length=3pt]}, thick},
    oarr/.style={-{Stealth[length=3pt]}, thick, dashed, gray},
    lbl/.style={font=\tiny, fill=white, inner sep=1pt}
]
\node[role] (I) at (0,2.4) {\textbf{ผู้ออกข้ออ้าง}};
\node[role] (S) at (-2.5,2.4) {\textbf{เป้าหมาย}};
\node[opt] (A) at (2.5,1.2) {\textbf{ผู้มีอำนาจ}};
\node[opt] (W) at (-2.5,0) {\textbf{พยาน}};
\node[role] (V) at (2.5,-0.6) {\textbf{ผู้ตรวจสอบ}};
\node[role, fill=black!5, draw=gray] (N) at (0,-1.8) {RSN};

\draw[arr] (I) -- node[lbl] {เกี่ยวกับ} (S);
\draw[oarr] (I) -- node[lbl,sloped] {รับรอง} (A);
\draw[oarr] (I) -- node[lbl,sloped] {ท้าทาย} (W);
\draw[arr] (I) -- node[lbl,sloped] {ร้องขอ} (V);
\draw[oarr] (W) -- node[lbl] {เฝ้าสังเกต} (V);
\draw[arr] (V) -- node[lbl] {เผยแพร่} (N);
\end{tikzpicture}
\caption{บทบาท DID ในธุรกรรมการตรวจสอบ เส้นประ = เป็นทางเลือก (ผู้มีอำนาจ พยาน)}
\label{fig:roles}
\end{figure}

\subsection{ความไม่อาจทุจริตได้: สี่แรง}

ความไม่อาจทุจริตได้ของ RSN มิได้เกิดจากกลไกเดียว แต่เกิดจากสี่แรงที่ทำงานพร้อมกัน ไม่มีแรงใดเพียงพอในตัวเอง มีเพียงการผสานกันของแรงเหล่านี้เท่านั้นที่ทำให้ความพยายามทุจริตไม่คุ้มค่าในเชิงเศรษฐกิจ

\textbf{เป้าหมายที่คาดเดาไม่ได้} ผู้ตรวจสอบสำหรับแต่ละข้ออ้างถูกคัดเลือกแบบไม่กำหนดตายตัวผ่านการแฮชแบบสม่ำเสมอ (หัวข้อ 5.3) ผู้โจมตีไม่สามารถล็อกเป้าผู้ตรวจสอบรายใดรายหนึ่งไว้ล่วงหน้าเพื่อติดสินบนได้ เพราะอัตลักษณ์ของผู้ตรวจสอบเป็นฟังก์ชันของเนื้อหาข้ออ้างและการทำซ้ำครั้งก่อนๆ มิใช่ตัวเลือกของผู้ออกข้ออ้าง

\textbf{แรงเหนี่ยวนำด้านชื่อเสียง---ขึ้นช้า ลงเร็ว} ชื่อเสียงได้มาได้ทีละน้อยเท่านั้น ผ่านพฤติกรรมที่สม่ำเสมออย่างยั่งยืน อย่างไรก็ตาม ชื่อเสียงอาจสูญเสียไปอย่างมหันต์ด้วยการรับรองที่ฉ้อฉลเพียงครั้งเดียว (หัวข้อ 6.2) ความไม่สมมาตรนี้หมายความว่าชื่อเสียงที่สั่งสมมาหลายปีอาจถูกทำลายด้วยการรับรองที่ไม่ซื่อสัตย์เพียงครั้งเดียว กลยุทธ์ที่เหมาะสมที่สุดจึงยังคงเป็นการปฏิเสธข้ออ้างที่น่าสงสัย

\textbf{คำมั่นสาธารณะ} ผู้ถือ DID ทุกคนประกาศนโยบายของตนต่อสาธารณะ---ชุดกฎที่เครื่องอ่านได้ซึ่งพวกเขาให้คำมั่นว่าจะปฏิบัติตาม ความคลาดเคลื่อนระหว่างคำประกาศกับพฤติกรรมที่แท้จริงสามารถตรวจจับได้และถูกตั้งราคาในฐานะความผิดด้านชื่อเสียงที่ร้ายแรงกว่าการละเมิดพื้นฐาน (หัวข้อ 7.3) ความหน้าไหว้หลังหลอกจึงมีต้นทุนแพงกว่าความไม่ซื่อสัตย์โดยตรง

\textbf{ความไม่สมมาตรของต้นทุนการโจมตีแบบตาข่าย} ต้นทุนของการเซ็นเซอร์หรือการทำให้เครือข่ายไร้เสถียรภาพเติบโตแบบไม่เป็นเชิงเส้นตามขนาดและความหนาแน่นของเครือข่าย ในขณะที่ต้นทุนของการยอมทนต่อมันคงที่ เพื่อให้การเซ็นเซอร์คุ้มค่า ผู้กระทำแบบรวมศูนย์จะต้องใช้มาตรการนั้นทั่วทั้งเครือข่าย---ซึ่งต้องอาศัยมาตรการที่ไม่ได้สัดส่วนกับผลประโยชน์ใดที่นึกออก (หัวข้อ 10)

% ============================================================
% 4. DECENTRALIZED IDENTITY MODEL
% ============================================================
\section{แบบจำลองอัตลักษณ์แบบกระจายศูนย์}

\subsection{DID พร้อมคุณสมบัติพิเศษ}

RSN ขยายข้อกำหนด W3C DID Core~\cite{did2022} ด้วย: \textbf{กฎที่ประกาศไว้} (คำมั่นด้านพฤติกรรมที่เครื่องอ่านได้), \textbf{ข้อมูลอภิพันธุ์ชื่อเสียง} (คะแนนพฤติกรรมที่รวบรวมแล้ว) และ \textbf{การกำหนดเส้นทางในเครือข่าย} (เกตเวย์แบบหัวหอมที่เปลี่ยนแปลงได้ แยกออกจากตำแหน่งบนวงแหวนที่คงที่)

\subsection{วงจรชีวิตของข้ออ้าง}

ข้ออ้างหนึ่งดำเนินไปผ่านหกขั้นตอน (รูปที่~\ref{fig:lifecycle}): การประกอบ การรับรองโดยผู้มีอำนาจ (ตามทางเลือก) การคัดเลือกผู้ตรวจสอบผ่านการแฮชแบบสม่ำเสมอ การตัดสินใจตรวจสอบ (ยอมรับหรือปฏิเสธพร้อมทำซ้ำ) การเผยแพร่ และการปรับปรุงชื่อเสียงสำหรับทุกฝ่าย

\begin{figure}[ht]
\centering
\begin{tikzpicture}[
    stp/.style={draw, rounded corners=2pt, minimum width=1.3cm, minimum height=0.45cm, align=center, font=\tiny, fill=black!5},
    arr/.style={-{Stealth[length=3pt]}, thick},
    lbl/.style={font=\tiny, fill=white, inner sep=1pt}
]
\node[stp] (comp) at (0,0) {ประกอบ\\ข้ออ้าง};
\node[stp, dashed] (auth) at (2.0,0) {ผู้มีอำนาจ\\รับรอง};
\node[stp] (hash) at (4.0,0) {แฮช \&\\แมปวงแหวน};
\node[stp] (ver) at (2.0,-1.6) {ผู้ตรวจสอบ\\ตรวจ};
\node[stp, double] (pub) at (0,-1.6) {เผยแพร่แล้ว};
\node[stp] (rej) at (4.0,-1.6) {ปฏิเสธ\\ต้นทุน $\uparrow$};

\draw[arr] (comp) -- (auth);
\draw[arr] (auth) -- (hash);
\draw[arr] (hash) -- (ver);
\draw[arr] (ver) -- node[lbl] {\checkmark} (pub);
\draw[arr] (ver) -- node[lbl] {$\times$} (rej);
\draw[arr] (rej) -- ++(0,0.8) -| (hash);
\end{tikzpicture}
\caption{วงจรชีวิตของข้ออ้างพร้อมวงวนการทำซ้ำ}
\label{fig:lifecycle}
\end{figure}

\subsection{ความสัมพันธ์กับ W3C DID Core}

แบบจำลองอัตลักษณ์ของ RSN เป็นเซตยิ่งยวด (superset) ที่ถูกต้องของข้อกำหนด W3C DID Core~\cite{did2022} DID ทั้งหมดของ RSN เป็น W3C DID ที่ถูกต้อง ส่วนขยายเฉพาะของ RSN ถูกเข้ารหัสเป็นคุณสมบัติของเอกสาร DID ที่นิยามไว้ในข้อกำหนดวิธีการ DID เฉพาะ ซึ่งเข้ากันได้กับโครงสร้างพื้นฐานที่มีอยู่เดิม เช่น ION~\cite{buchner2021} และ Ceramic~\cite{ceramic2023}

% ============================================================
% 5. VERIFICATION AND PROPAGATION
% ============================================================
\section{การตรวจสอบและการแพร่กระจายข้อมูล}

\subsection{ต้นทุนของการป้อนข้อมูลเข้าสู่ระบบ}

หลักการทางเศรษฐกิจแกนกลางของ RSN คือความไม่สมมาตรระหว่างการอ่านและการเขียน การอ่านข้อมูลชื่อเสียงเข้าใกล้ต้นทุนส่วนเพิ่มเป็นศูนย์สำหรับผู้เข้าร่วมที่เป็นส่วนหนึ่งของเครือข่าย DID และมีสิทธิ์เข้าถึงตามสัดส่วนของชื่อเสียงของตน---ผู้มาใหม่ต้องค่อยๆ หาสิทธิ์เข้าถึงที่กว้างขึ้น (ดูการต้านการเอาเปรียบ) การเขียน---เข้าใจในฐานะการทำให้ชื่อเสียงก่อรูปเป็นสสารในเครือข่ายอย่างถาวร มิใช่การกระทำทางเทคนิคเพียงครั้งเดียว---ต้องอาศัยการลงทุนสะสมที่พิสูจน์ได้ในสามมิติ ได้แก่ \textbf{เวลา} (หลายปีของชีวิตที่ทุ่มเทให้พฤติกรรมที่สม่ำเสมอ คำมั่นที่ทำได้จริง และความสัมพันธ์ที่บ่มเพาะ), \textbf{พลังงาน} (ความพยายามที่ลงทุนไปกับการปฏิบัติต่อกันอย่างซื่อสัตย์ ความใส่ใจในหุ้นส่วน และความน่าไว้วางใจในระยะยาว) และ \textbf{เงิน} (การลงทุนในชื่อเสียงของตนเอง---การพัฒนาวิชาชีพ คุณภาพของงาน การชดใช้ความเสียหายที่ก่อขึ้น) ชื่อเสียงไม่สามารถซื้อได้ในทันที---มันคือผลลัพธ์ของการสั่งสมตลอดชีวิตซึ่งระบบเพียงทำให้มองเห็นได้และโอนย้ายข้ามบริบทได้ ต้นทุนทางเทคนิคครั้งเดียวของโปรโตคอล (ลายเซ็นทางวิทยาการเข้ารหัส ค่าธรรมเนียมผู้ตรวจสอบ) นั้นเล็กน้อยเมื่อเทียบกับการลงทุนพื้นฐานนี้

\subsection{กราฟสังคมและความลึกของการติดต่อ}

RSN โดยพื้นฐานแล้วคือเครือข่ายสังคม แต่ละ DID เพิ่มผู้ติดต่อ (DID อื่นๆ) ที่อนุญาตการเชื่อมต่อ ผู้ติดต่อเหล่านี้มีผู้ติดต่อของตนเอง และต่อเนื่องไปเช่นนี้ การค้นหาผู้ตรวจสอบด้วยอัลกอริทึมดำเนินการภายในความลึก $h$ ของผู้ติดต่อที่เชื่อมโยงกันซึ่งกำหนดค่าได้ (เช่น $h=3$: ผู้ติดต่อโดยตรงของตน ผู้ติดต่อของพวกเขา และผู้ติดต่อของผู้ติดต่อ) สถาปัตยกรรมนี้ไม่ต้องการบล็อกเชนระดับโลกเพียงหนึ่งเดียว---แต่เอื้อให้เกิดชุมชน/กลุ่มก้อนโดยธรรมชาติ ที่มีการทับซ้อนเข้าไปในชุมชนอื่นๆ ผู้เข้าร่วม DID ทุกคนต้องมี\textbf{นโยบาย}ที่เผยแพร่ไว้---ชุดกฎที่เครื่องอ่านได้ซึ่งกำกับว่าคำร้องขอที่เข้ามาจะได้รับการประเมินอย่างไร การละเมิดนโยบายจะถูกบันทึกในเครือข่ายในฐานะสิ่งประดิษฐ์เชิงลบ (negative artifacts)

\subsection{การคัดเลือกผู้ตรวจสอบด้วยอัลกอริทึม}

โปรโตคอลการตรวจสอบใช้การแฮชแบบสม่ำเสมอ~\cite{karger1997} (รูปที่~\ref{fig:verifier}) การตรวจสอบเป็นบริการที่มีค่าใช้จ่าย ผู้ตรวจสอบดำเนินการโดยมีค่าธรรมเนียม ก่อให้เกิดตลาดที่การแข่งขันผลักดันการตั้งราคา ความเร็ว และความน่าเชื่อถือ

\textbf{ขั้นที่ 1} เอกสารข้ออ้างถูกต่อเข้ากับผลลัพธ์แฮชของการทำซ้ำครั้งก่อน (หรือ $\varnothing$ ในการทำซ้ำครั้งแรก) แล้วนำไปแฮช:
\begin{equation}
\text{pos} = f_h(\text{claim} \| \text{prev\_hash})
\label{eq:hash}
\end{equation}

อัลกอริทึมต้องรวม\emph{เส้นทางที่สมบูรณ์}ของคำร้องขอและการตอบสนองก่อนหน้าทั้งหมด---มิใช่เพียงแฮชของการทำซ้ำครั้งก่อน การตอบสนองฉบับเต็มของผู้ตรวจสอบแต่ละราย (การยอมรับ การปฏิเสธ ราคาที่เสนอ เหตุผล) ถูกต่อเข้ากับเอกสารที่เติบโตขึ้น ซึ่งตรวจสอบได้ผ่านลายเซ็นทางวิทยาการเข้ารหัสในทุกขั้นตอน

\textbf{ขั้นที่ 2} แฮชถูกแมปไปยัง $N$ ตำแหน่งบนวงแหวนแฮชแบบสม่ำเสมอ กระจายอย่างสม่ำเสมอ (เช่น สำหรับ $N=4$: $+0^{\circ}, +90^{\circ}, +180^{\circ}, +270^{\circ}$) จากแต่ละตำแหน่ง การค้นหาตามเข็มนาฬิกาจะเลือก DID ที่ใกล้ที่สุดเป็นผู้ตรวจสอบตัวเลือก (รูปที่~\ref{fig:multipoint}) $N$ เป็นพารามิเตอร์แปรผันที่อาจขึ้นอยู่กับเปอร์เซ็นต์ของ DID ที่กำลังใช้งานอยู่ในปัจจุบัน เวลาของวัน หรือการตอบสนองในอดีตของเครือข่าย

\textbf{ขั้นที่ 3} ผู้ตรวจสอบยอมรับ (เผยแพร่ พร้อมเดิมพันชื่อเสียง) หรือปฏิเสธ (กลับไปยังขั้นที่ 1 พร้อมแฮชของการปฏิเสธ) เมื่อโหนดหนึ่งไม่ตอบสนองทั้งที่ประกาศสถานะออนไลน์ คำร้องขอครั้งถัดไปต้องรวมการตอบสนองทั้งหมดจากผู้ตรวจสอบตัวเลือกทั้ง $N$ รายก่อนหน้า

\textbf{การต้านการเอาเปรียบ (anti-leeching)} DID ที่เป็นเป้าหมายอาจตั้งค่าธรรมเนียมหรือข้อจำกัดการเข้าถึงสำหรับคำสอบถามจาก DID ใหม่ที่มีชื่อเสียงน้อย กลไกแบบเป็นลำดับขั้นนี้ (มิใช่แบบทวิภาค) บังคับให้ผู้มาใหม่ต้องสร้างชื่อเสียงผ่านกิจกรรมที่แท้จริงก่อนที่จะบริโภคข้อมูลของผู้อื่นได้อย่างเสรี

\begin{figure}[ht]
\centering
\begin{tikzpicture}[
    box/.style={draw, rounded corners=2pt, minimum width=1.3cm, minimum height=0.45cm, align=center, font=\scriptsize, fill=black!5},
    dec/.style={draw, diamond, aspect=2.2, minimum width=0.9cm, align=center, font=\scriptsize, fill=black!5},
    arr/.style={-{Stealth[length=3pt]}, thick},
    lbl/.style={font=\tiny, fill=white, inner sep=1pt}
]
\node[box] (iss) at (0,0) {ผู้ออกข้ออ้าง};
\node[box] (hash) at (0,-1.1) {$f_h$(claim $\|$\\prev\_hash)};
\node[box] (ring) at (0,-2.2) {แมปวงแหวน\\ตามเข็ม};
\node[box, dashed] (spam) at (0,-3.2) {เงินมัดจำ\\ต้านสแปม};
\node[dec] (check) at (0,-4.4) {ตรวจ\\นโยบาย};
\node[box] (rej) at (2,-5.6) {ทำซ้ำ\\ครั้งถัดไป};
\node[box] (fee) at (0,-5.6) {ค่าธรรมเนียมสุดท้าย =\\$f$(ข้อมูล, ชื่อเสียง, นโยบาย)};
\node[box, double] (pub) at (0,-6.8) {เผยแพร่แล้ว};

\draw[arr] (iss) -- (hash);
\draw[arr] (hash) -- (ring);
\draw[arr] (ring) -- (spam);
\draw[arr] (spam) -- (check);
\draw[arr] (check) -- node[lbl] {$\times$} (rej);
\draw[arr] (check) -- node[lbl] {\checkmark} (fee);
\draw[arr] (fee) -- (pub);
\draw[arr] (rej.east) -- ++(0.5,0) |- (hash.east);
\end{tikzpicture}
\caption{โปรโตคอลการคัดเลือกผู้ตรวจสอบ เงินมัดจำต้านสแปมเป็นทางเลือกและอาจคืนได้ ค่าธรรมเนียมสุดท้ายชำระเมื่อได้รับการยอมรับเท่านั้น}
\label{fig:verifier}
\end{figure}

การปฏิเสธแต่ละครั้งจะต่อการตอบสนองฉบับเต็มของผู้ตรวจสอบ (รวมทั้งเหตุผลและเงื่อนไข) เข้ากับเอกสารข้ออ้าง ทำให้เนื้อหาขยายออก จากเอกสารที่ขยายแล้ว แฮชใหม่จะถูกคำนวณแบบไม่กำหนดตายตัว แมปไปยังตำแหน่งบนวงแหวนที่ต่างออกไป และคัดเลือกผู้ตรวจสอบที่ต่างออกไป การจัดทำเอกสารเส้นทางฉบับเต็มเป็นข้อบังคับ---ผู้ออกข้ออ้างไม่สามารถคาดเดาหรือมีอิทธิพลต่อผู้ตรวจสอบรายถัดไป หากผู้ตรวจสอบเพิกเฉยต่อคำร้องขอทั้งที่มีนโยบายที่ประกาศไว้เข้ากันได้ พยาน (หากมี) จะบันทึกเหตุการณ์นี้ และผู้ออกข้ออ้างอาจแนบหลักฐานจากพยานได้เมื่อเผยแพร่ครั้งสุดท้าย

\begin{figure}[ht]
\centering
\begin{tikzpicture}[scale=0.65]
% Ring
\draw[thick] (0,0) circle (2.2cm);
% DID nodes on ring
\foreach \a in {10,55,80,120,150,195,230,260,305,340} {
  \fill[black!40] (\a:2.2) circle (1.5pt);
}
% Hash offset arrow pointing to initial position
\draw[-{Stealth[length=3pt]}, thick, black!60] (0,0) -- node[font=\tiny,above,sloped,pos=0.6] {$f_h$} (37:2.2);
\fill[black] (37:2.2) circle (2pt);
\node[font=\tiny,anchor=south west] at (37:2.35) {$h_0$};
% N=4 points offset from h0 by +0, +90, +180, +270
\foreach \offset/\l in {0/P1,90/P2,180/P3,270/P4} {
  \pgfmathsetmacro{\ang}{37+\offset}
  \node[fill=black, circle, inner sep=1.5pt] (p\offset) at (\ang:2.2) {};
  \node[font=\tiny,font=\bfseries] at (\ang:2.75) {\l};
  % clockwise search arc
  \draw[-{Stealth[length=2.5pt]}, thick, gray] (\ang:2.2) arc (\ang:\ang+30:2.2);
}
\node[font=\scriptsize, align=center] at (0,0) {วงแหวน\\แฮช};
\node[font=\tiny, align=center] at (0,-3.2) {$h_0 = f_h(\text{claim})$, แล้ว $+0^{\circ},+90^{\circ},+180^{\circ},+270^{\circ}$\\แต่ละจุด $\rightarrow$ ค้นหาตามเข็มนาฬิกา};
\end{tikzpicture}
\caption{การคัดเลือกแบบหลายจุด ($N{=}4$) แฮช $h_0$ กำหนดออฟเซต; 4 จุดกระจายอย่างสม่ำเสมอจาก $h_0$}
\label{fig:multipoint}
\end{figure}

\subsection{โปรโตคอลพยาน}

บทบาท\textbf{พยาน}ให้ความรับผิดชอบแบบนิรนามต่อความซื่อสัตย์ของผู้ตรวจสอบ ผ่านโปรโตคอลรหัสท้าทายทางวิทยาการเข้ารหัส:

\textbf{ขั้น W1} ผู้ร้องขอเซ็นคำร้องขอการตรวจสอบและส่งไปยังพยาน $N_w$ ราย---ผู้ติดต่อจากเครือข่ายสังคมของผู้ร้องขอ หรือผู้ให้บริการพยานเฉพาะทางที่ดำเนินการโดยมีค่าธรรมเนียม

\textbf{ขั้น W2} พยาน $w_i$ แต่ละรายเซ็นคำร้องขอที่เซ็นแล้วทั้งฉบับ เก็บลายเซ็นต้นฉบับ $\sigma_i$ ไว้ และคำนวณรหัสท้าทาย $c_i = h(\sigma_i)$ รหัสท้าทายจะถูกต่อเข้ากับคำร้องขอ

\textbf{ขั้น W3} คำร้องขอซึ่งบัดนี้พกพารหัสท้าทาย $N_w$ รหัส ถูกส่งไปยังผู้ตรวจสอบ ผู้ตรวจสอบมองเห็นรหัสเหล่านั้นแต่\emph{ไม่สามารถระบุได้}ว่าใครคือพยาน หรือรหัสเหล่านั้นตรงกับลายเซ็นจริงหรือไม่

\textbf{ขั้น W4 (ผู้ตรวจสอบซื่อสัตย์)} หากผู้ตรวจสอบตอบสนองตามนโยบายที่ประกาศไว้ รหัสท้าทายจะคงความทึบแสงไว้ พยานจะไม่ถูกเปิดเผยเลย ไม่มีการเผยแพร่ข้อมูลเพิ่มเติม

\textbf{ขั้น W5 (การละเมิดนโยบาย)} หากผู้ตรวจสอบละเมิดนโยบายของตน (เพิกเฉยต่อคำร้องขอ ตอบสนองอย่างไม่สอดคล้อง) ผู้ร้องขอจะถือลายเซ็นต้นฉบับ $\sigma_i$ ของพยานไว้ ลายเซ็นเหล่านี้สามารถเผยแพร่เป็นข้อมูลประกอบได้ ใครก็ตามสามารถตรวจสอบ $c_i = h(\sigma_i)$ พิสูจน์ได้ว่าพยานอยู่ในเหตุการณ์ ผู้ร้องขอสามารถเผยแพร่ได้โดยอิสระ---ผู้ตรวจสอบได้เซ็นรหัสท้าทายไว้ในการตอบสนองของตนแล้ว

\textbf{คุณสมบัติการหลอกลวง (bluff property)} ผู้ร้องขออาจแทนที่รหัสท้าทายบางส่วนหรือทั้งหมดด้วยค่าสุ่ม ผู้ตรวจสอบไม่สามารถแยกแยะรหัสของจริง (ที่หนุนหลังโดยผู้มีอำนาจชื่อเสียงสูง) ออกจากสัญญาณรบกวนได้ \emph{ความไม่แน่นอน}นี้คือกลไกการบังคับใช้ ทุกคำร้องขอพกพาความเสี่ยงที่ว่าผู้มีอำนาจที่ได้รับการเคารพกำลังเฝ้าดูอยู่แบบนิรนาม ต้นทุนของการรักษาแรงกดดันนี้เข้าใกล้ศูนย์ (การสุ่มตัวเลข) ในขณะที่ต้นทุนที่อาจเกิดขึ้นจากความไม่ซื่อสัตย์ของผู้ตรวจสอบนั้นมหันต์ พฤติกรรมซื่อสัตย์จึงถูกจูงใจแม้ในภาวะที่ไม่มีพยานอยู่จริง

\textbf{ผู้มีอำนาจในฐานะพยาน} ผู้มีอำนาจที่รับรองข้ออ้างหนึ่งอาจรวมบริการพยานเข้าไปด้วย: ``ข้าพเจ้ารับรองข้ออ้างของท่านและทำหน้าที่เป็นพยานนิรนามระหว่างการตรวจสอบ'' นี่คือโมเดลธุรกิจที่เป็นธรรมชาติ---ผู้มีอำนาจมีบริบทอยู่แล้ว อย่างไรก็ตาม บทบาททั้งสองเป็นอิสระต่อกันในเชิงตรรกะ

\subsection{หลักการความสุดโต่งราคาแพง}

ช่องทางต้นทุนสามช่องทางทบต้นกันพร้อมกัน (รูปที่~\ref{fig:radicalism}):

\textbf{ช่องทางที่ 1: ต้นทุนการทำซ้ำ} การปฏิเสธแต่ละครั้งบังคับให้เกิดการทำซ้ำใหม่ที่กินเวลาและทรัพยากร เติบโตแบบเชิงเส้น

\textbf{ช่องทางที่ 2: การพองตัวของเอกสาร} การทำซ้ำแต่ละครั้งเพิ่มการตอบสนองฉบับเต็มของผู้ตรวจสอบเข้าไปในเอกสารข้ออ้าง การเติบโตเป็นเชิงเส้น แต่มีค่าฐานเริ่มต้น นั่นคือ เพย์โหลดเริ่มต้น (เนื้อหาข้ออ้าง การรับรองของผู้มีอำนาจ ลายเซ็นทางวิทยาการเข้ารหัส) มีขนาดที่ไม่ใช่เรื่องเล็กน้อยอยู่แล้วที่ $B_0$ ขนาดรวมหลังจากทำซ้ำ $k$ ครั้ง: $S(k) = B_0 + k \cdot \Delta$ โดยที่ $\Delta$ คือขนาดเฉลี่ยของการตอบสนอง

\textbf{ช่องทางที่ 3: ส่วนเพิ่มความเสี่ยงของผู้ตรวจสอบ} ความเสี่ยงเป็นเรื่องอัตวิสัย ผู้ตรวจสอบประเมินว่านโยบายที่ประกาศไว้ของ DID ผู้ร้องขอขัดแย้งกับผลประโยชน์ทางการค้า สังคม หรือการเมืองของผู้ตรวจสอบเองหรือไม่ ส่วนเพิ่มอาจทวีขึ้นแบบเลขชี้กำลังตามระยะห่างที่รับรู้จาก ``อุดมคติ'' ของผู้ตรวจสอบ: $P(d) = \alpha \cdot e^{\beta d}$ โดยที่ $d$ คือมาตรระยะห่างเชิงอัตวิสัยของผู้ตรวจสอบ เมื่อเลยเส้นแบ่งเขตห้ามเข้าส่วนตัวไปแล้ว ผู้ตรวจสอบอาจปฏิเสธโดยสิ้นเชิง

\begin{figure}[ht]
\centering
\begin{tikzpicture}
\begin{axis}[
    width=\columnwidth,
    height=4.5cm,
    xlabel={\footnotesize ความสุดโต่ง},
    ylabel={\footnotesize ต้นทุนสัมพัทธ์ (log)},
    ymode=log,
    xmin=0, xmax=100,
    ymin=1, ymax=10000,
    grid=major,
    grid style={gray!20},
    legend style={at={(0.03,0.97)},anchor=north west,font=\tiny,draw=gray!50},
    tick label style={font=\tiny},
    label style={font=\footnotesize},
]
\addplot[black, thick, domain=0:100, samples=50] {1 + 0.3*x};
\addlegendentry{ต้นทุนการทำซ้ำ (เชิงเส้น)}
\addplot[black, thick, dashed, domain=0:100, samples=50] {1 + 0.15*x};
\addlegendentry{การพองตัวของเอกสาร (เชิงเส้น)}
\addplot[black, thick, dotted, line width=1.2pt, domain=0:100, samples=100] {exp(0.08*x)};
\addlegendentry{ส่วนเพิ่มความเสี่ยง (เลขชี้กำลัง)}
\end{axis}
\end{tikzpicture}
\caption{การทวีขึ้นของต้นทุนใน 3 ช่องทาง ส่วนเพิ่มความเสี่ยงครอบงำที่ระดับความสุดโต่งสูง}
\label{fig:radicalism}
\end{figure}

ที่สำคัญ นี่ไม่ใช่การเซ็นเซอร์ ไม่มีข้ออ้างใดถูกห้าม ระบบตั้งราคาความเสี่ยง (ตารางที่~\ref{tab:costs})

\begin{table}[ht]
\centering
\caption{การเปรียบเทียบต้นทุนตามประเภทข้ออ้าง}
\label{tab:costs}
\footnotesize
\begin{tabular}{@{}lcccc@{}}
\toprule
\textbf{ประเภท} & \textbf{ทำซ้ำ} & \textbf{เอกสาร} & \textbf{ค่าธรรม.} & \textbf{รวม} \\
\midrule
น่าเชื่อถือ & $k_{\min}$ & $B_0$ & $P_{\min}$ & ต่ำ \\
เป็นที่ถกเถียง & $k_{\text{med}}$ & $B_0{+}k\Delta$ & $\alpha e^{\beta d}$ & ปานกลาง \\
สุดโต่ง & $k_{\max}$ & $\gg B_0$ & $\gg P_{\min}$ & สูงมาก \\
ฉ้อฉล & $\to\infty$ & --- & --- & เผยแพร่ไม่ได้ \\
\bottomrule
\end{tabular}
\end{table}

% ============================================================
% 6. REPUTATION AND RISK
% ============================================================
\section{ชื่อเสียงและการประเมินความเสี่ยง}

\subsection{แบบจำลองการสั่งสมชื่อเสียง}

ชื่อเสียงแสดงคุณสมบัติสามประการ ได้แก่ \textbf{การสั่งสมอย่างเชื่องช้า} (การเติบโตทีละน้อยผ่านพฤติกรรมที่สม่ำเสมอ), \textbf{การถูกทำลายอย่างรวดเร็ว} (การฉ้อฉลเพียงครั้งเดียวสามารถลดคะแนนได้อย่างมหันต์) และ \textbf{การประเมินเป็นรายบุคคล} (แต่ละฝ่ายที่บริโภคข้อมูลอาจใช้ฟังก์ชันการรวบรวมของตนเอง)

\subsection{ผู้มีอำนาจที่น่าเชื่อถือ}

ผู้มีอำนาจที่น่าเชื่อถือทบทวนข้ออ้าง และหากพึงพอใจ ก็ร่วมเซ็นรับรอง---โดยเดิมพันชื่อเสียงของตนเอง (รูปที่~\ref{fig:authority}) ความไม่สมมาตรนี้เป็นไปโดยการออกแบบ นั่นคือ การได้มาซึ่งชื่อเสียงเป็นไปอย่างเชื่องช้า (เพิ่มขึ้น $+\delta$ ต่อการรับรองที่ซื่อสัตย์แต่ละครั้ง) ในขณะที่การสูญเสียมันเป็นไปอย่างมหันต์ ($-\Delta \gg \delta$ ต่อการรับรองอันเป็นเท็จแต่ละครั้ง; เพื่อเป็นตัวอย่าง เช่น $+2\%$ เทียบกับ $-70\%$) กลยุทธ์ที่เหมาะสมที่สุดคือการปฏิเสธข้ออ้างที่น่าสงสัยเสมอ ค่าเฉพาะของ $\delta$ และ $\Delta$ เป็นพารามิเตอร์ของระบบ

\begin{figure}[ht]
\centering
\begin{tikzpicture}[
    box/.style={draw, rounded corners=2pt, minimum width=2cm, minimum height=0.6cm, align=center, font=\scriptsize, fill=black!5},
    arr/.style={-{Stealth[length=4pt]}, thick}
]
\node[box] (exam) at (0,0) {ผู้มีอำนาจตรวจสอบ\\ชื่อเสียง: $R$};
\node[box] (true) at (-2,-1.3) {จริง: $R{\to}R{+}\delta$\\ลูกค้ามากขึ้น};
\node[box] (false) at (2,-1.3) {เท็จ: $R{\to}R{-}\Delta$\\ลูกค้าจากไป};
\node[box] (insight) at (0,-2.6) {ได้: ช้า ($+\delta$)\\เสีย: ทันที ($-\Delta$)};

\draw[arr] (exam) -- node[left,font=\tiny] {ความจริง} (true);
\draw[arr] (exam) -- node[right,font=\tiny] {การโกหก} (false);
\draw[arr] (true) -- (insight);
\draw[arr] (false) -- (insight);
\end{tikzpicture}
\caption{ผู้มีอำนาจที่น่าเชื่อถือ---แรงจูงใจแบบไม่สมมาตร}
\label{fig:authority}
\end{figure}

\subsection{ชื่อเสียงหลายมิติ}

ชื่อเสียงไม่ใช่ปริมาณสเกลาร์ แต่เป็นเวกเตอร์ที่ทอดข้ามหลายมิติ ได้แก่ ความน่าเชื่อถือทางการเงิน ความซื่อสัตย์ส่วนบุคคล ความสามารถทางวิชาชีพ การมีส่วนร่วมทางพลเมือง และอื่นๆ (รูปที่~\ref{fig:reputation-radar}) ผู้ให้บริการประเมินที่แตกต่างกันฉายเวกเตอร์นี้แตกต่างกันขึ้นอยู่กับบริบท---ผู้ให้กู้ให้ความสำคัญกับความน่าเชื่อถือทางการเงิน นายจ้างให้ความสำคัญกับความสามารถทางวิชาชีพ เพื่อนบ้านให้ความสำคัญกับพฤติกรรมทางพลเมือง ไม่มีคะแนนชื่อเสียงที่ ``ถูกต้อง'' เพียงหนึ่งเดียว มีเพียงมุมมองที่ต่างกัน

\begin{figure}[ht]
\centering
\begin{tikzpicture}[scale=0.55]
\foreach \a/\l in {90/การเงิน,162/ความซื่อสัตย์,234/วิชาชีพ,306/พลเมือง,18/สังคม} {
  \draw[gray!40] (0,0) -- (\a:2.5);
  \node[font=\tiny, align=center] at (\a:3.1) {\l};
  \foreach \r in {0.8,1.6,2.4} {
    \fill[black!15] (\a:\r) circle (0.5pt);
  }
}
\draw[thick, fill=black!10] (90:2.0) -- (162:1.2) -- (234:1.8) -- (306:0.9) -- (18:1.5) -- cycle;
\draw[thick, dashed] (90:1.4) -- (162:2.1) -- (234:0.8) -- (306:2.0) -- (18:1.1) -- cycle;
\node[font=\tiny] at (0,-3.3) {เส้นทึบ: DID A \quad เส้นประ: DID B};
\end{tikzpicture}
\caption{เวกเตอร์ชื่อเสียงหลายมิติ DID ที่ต่างกันแสดงโปรไฟล์ที่ต่างกันในแต่ละมิติ}
\label{fig:reputation-radar}
\end{figure}

\subsection{การประเมินความเสี่ยงที่เป็นประชาธิปไตย}

RSN ทำให้การประเมินความเสี่ยงอย่างเป็นระบบเป็นประชาธิปไตย---ความสามารถที่ปัจจุบันกระจุกตัวอยู่ในสถาบันการเงิน แต่ประยุกต์ใช้ได้ไกลเกินกว่าแวดวงการธนาคาร กลุ่มบริษัทอุตสาหกรรมที่ประเมินความน่าเชื่อถือของซัพพลายเออร์ บริษัทประกันที่ตั้งราคาความเสี่ยงด้านพฤติกรรม การตรวจสอบสถานะกิจการสำหรับการควบรวมและซื้อกิจการ การประมาณอายุการใช้งานของอุปกรณ์ในภาคการผลิต---ทุกที่ที่ความเสี่ยงของคู่สัญญาต้องได้รับการประเมิน RSN มอบทางเลือกที่กระจายศูนย์และขับเคลื่อนด้วยตลาด ตลาดของบริการตีความจะแปลข้อมูลชื่อเสียงหลายมิติดิบให้เป็นบทสรุปที่นำไปปฏิบัติได้ ทำให้การประเมินคู่สัญญาเข้าถึงได้สำหรับผู้เข้าร่วมทุกคนที่ต้นทุนส่วนเพิ่ม

% ============================================================
% 7. CONSENSUS AND DISPUTE RESOLUTION
% ============================================================
\section{ฉันทามติและการระงับข้อพิพาท}

\subsection{แบบจำลองความยุติธรรมแบบกระจายศูนย์}

RSN วางกรอบความยุติธรรมใหม่ในฐานะปัญหาการเจรจาต่อรองแบบสองทาง ภัยคุกคามของความเสียหายต่อชื่อเสียงที่ถาวรและตรวจสอบได้เป็นเครื่องยับยั้งที่มีประสิทธิภาพมากกว่ากระบวนการทางกฎหมายที่ฝ่ายผู้เสียหายไม่มีปัญญาจ่าย

\subsection{ฉันทามติแบบผุดขึ้นเอง}

ผู้เข้าร่วมแต่ละคนรักษาเกณฑ์ความพึงพอใจส่วนบุคคลไว้ ฉันทามติโดยรวมของเครือข่ายผุดขึ้นในฐานะค่าเฉลี่ยทางสถิติของการเจรจาต่อรองแบบสองทางนับพันครั้ง (รูปที่~\ref{fig:consensus}) การตอบสนองที่สูงกว่าฉันทามติมาก (ป่าเถื่อน) มีราคาแพงในการเผยแพร่ การตอบสนองที่ใกล้ฉันทามติ (ได้สัดส่วน) มีราคาถูก ฉันทามติไม่ใช่กฎ---มันคือสัญญาณราคา

\begin{figure}[ht]
\centering
\begin{tikzpicture}[
    person/.style={draw, rounded corners=2pt, minimum width=1.5cm, minimum height=0.5cm, align=center, font=\scriptsize, fill=black!5},
    arr/.style={-{Stealth[length=4pt]}, thick}
]
\node[person] (a) at (-2,0) {เข้มงวด\\(สูง)};
\node[person] (b) at (0,0) {ปฏิบัตินิยม\\(กลาง)};
\node[person] (c) at (2,0) {ให้อภัย\\(ต่ำ)};
\node[person, fill=black!10, minimum width=3cm] (cons) at (0,-1.2) {ฉันทามติเครือข่าย\\= ค่าเฉลี่ยสถิติ};
\node[person] (exp) at (-2,-2.4) {ป่าเถื่อน\\ราคาแพง};
\node[person, double] (prop) at (0,-2.4) {ได้สัดส่วน\\ราคาถูก};
\node[person] (neg) at (2,-2.4) {ละเลย\\ราคาแพง};

\draw[arr] (a) -- (cons);
\draw[arr] (b) -- (cons);
\draw[arr] (c) -- (cons);
\draw[arr] (cons) -- (exp);
\draw[arr] (cons) -- (prop);
\draw[arr] (cons) -- (neg);
\end{tikzpicture}
\caption{ฉันทามติแบบผุดขึ้นเองจากเกณฑ์ความพึงพอใจของแต่ละบุคคล}
\label{fig:consensus}
\end{figure}

\subsection{การปรับเทียบบทลงโทษ}

ผู้ถือ DID แต่ละคนประกาศต่อสาธารณะถึงการตอบสนองต่อประเภทของการประพฤติผิดที่เฉพาะเจาะจง คำประกาศที่รุนแรงหรือผ่อนปรนอย่างไม่ได้สัดส่วนจะดึงดูดสัญญาณชื่อเสียงเชิงลบ คำประกาศที่ได้สัดส่วนจะได้รับการยอมรับโดยไม่มีแรงเสียดทาน---ระบบบทลงโทษที่ปรับเทียบตัวเอง

คำประกาศฉันทามติเองก็อยู่ภายใต้การตรวจจับความหน้าไหว้หลังหลอก การให้คำมั่นในเชิงประกาศต่อจุดยืนที่เป็นฉันทามติในขณะที่ประพฤติตนไม่สอดคล้องกัน ถือเป็นความผิดด้านชื่อเสียงที่ร้ายแรงกว่าการเบี่ยงเบนพื้นฐาน เพราะแสดงให้เห็นถึงการหลอกลวงโดยเจตนา

\subsection{การมอบหมาย}

กิจกรรมทั้งหมดภายใน DID หนึ่ง---โดยมีข้อยกเว้นหนึ่งประการ---สามารถมอบหมายให้ DID อื่นได้ \textbf{บทบาทเป้าหมาย---DID ที่ข้ออ้างกล่าวถึงพฤติกรรมของมัน---ไม่สามารถมอบหมายได้; มันถูกผูกติดอย่างไม่อาจโอนย้ายกับผู้ถืออัตลักษณ์ที่ข้ออ้างนั้นอ้างถึง} บทบาทที่ทำงานอยู่อื่นๆ---ผู้ออกข้ออ้าง ผู้มีอำนาจ พยาน ผู้ตรวจสอบ---สามารถโอนย้ายให้ DID ตัวแทนที่กำหนดได้ ไม่ว่าจะเพื่อการตรวจสอบ การส่งข้ออ้าง การมีส่วนร่วมในฉันทามติ หรือการระงับข้อพิพาท การมอบหมายสามารถเพิกถอนได้ทุกเมื่อ สิ่งนี้เปิดโอกาสให้เกิดความเชี่ยวชาญเฉพาะทาง (เช่น บริการตรวจสอบระดับมืออาชีพ) ในขณะที่ผู้ถือยังคงมีการควบคุมและความรับผิดชอบสูงสุด การมอบหมายใช้ได้ทั่วทั้งระบบและจำเป็นต่อการขยายขนาด

\subsection{จากศีลธรรมของปัจเจกสู่สัญญาประชาคมแบบผุดขึ้นเอง}

นโยบายที่ประกาศไว้ของแต่ละ DID เป็นการทำให้ศีลธรรมของปัจเจกเป็นรูปธรรม นั่นคือ สิ่งที่ผู้ถือถือว่ายอมรับได้ วิธีที่พวกเขาตอบสนองต่อพฤติกรรมเฉพาะ สิ่งที่พวกเขาเรียกร้องจากคู่สัญญา นโยบายเหล่านี้ไม่ได้ถูกกำหนดโดยผู้ออกแบบระบบ---แต่ถูกเลือกโดยผู้เข้าร่วมแต่ละคนและแสดงออกในรูปแบบที่เครื่องอ่านได้

การทำให้เป็นรูปธรรมนี้เปิดทางให้เกิดการผุดขึ้นสามขั้น ขั้นแรก ศีลธรรมของปัจเจกถูกเข้ารหัสเป็น\textbf{นโยบาย}---ชุดกฎที่ประกาศไว้ เกณฑ์ และฟังก์ชันการตอบสนองที่ DID ให้คำมั่นว่าจะปฏิบัติตาม ขั้นที่สอง ผลรวมของนโยบายทั้งหมดทั่วเครือข่ายก่อให้เกิด\textbf{จริยธรรมแบบผุดขึ้นเอง}---บรรทัดฐานที่สังเกตได้ทางสถิติซึ่งไม่มีผู้เข้าร่วมคนใดออกแบบไว้ แต่เกิดขึ้นจากปฏิสัมพันธ์ของจุดยืนทางศีลธรรมของปัจเจกนับพัน~\cite{durkheim1893} ขั้นที่สาม จริยธรรมแบบผุดขึ้นเองเหล่านี้ ซึ่งแสดงออกในฐานะบรรทัดฐานพฤติกรรมร่วมที่บังคับใช้ผ่านสัญญาณราคาแบบสองทาง ประกอบกันเป็น\textbf{สัญญาประชาคมที่วัดได้}---มิใช่โครงสร้างเชิงทฤษฎีที่ไม่มีใครเซ็น~\cite{rawls1971} แต่เป็นชุดกฎที่สังเกตได้เชิงประจักษ์ซึ่งหนุนหลังด้วยพฤติกรรมจริง

กระบวนการนี้สะท้อนข้อค้นพบจากทฤษฎีรากฐานทางศีลธรรม (moral foundations theory)~\cite{haidt2012}: ศีลธรรมของมนุษย์เป็นสิ่งที่รู้ได้โดยสัญชาตญาณ มีพหุนิยม และผันแปรตามวัฒนธรรม RSN ไม่ต้องการฉันทามติทางศีลธรรมเป็นปัจจัยนำเข้า แต่ผลิตฉันทามติเชิงพฤติกรรมเป็นผลลัพธ์ สัญญาประชาคมที่เกิดขึ้นไม่ใช่สิ่งสถิต---มันวิวัฒน์ไปตามที่ผู้เข้าร่วมเข้ามา ออกไป และปรับนโยบายของตนเพื่อตอบสนองต่อการป้อนกลับของเครือข่าย ต่างจากทฤษฎีสัญญาประชาคมแบบดั้งเดิม สัญญานี้เพิกถอนได้ สังเกตได้ และบังคับใช้ได้โดยไม่ต้องมีองค์อธิปัตย์

% ============================================================
% 8. ECONOMIC MODEL
% ============================================================
\section{แบบจำลองทางเศรษฐกิจ}

หัวข้อก่อนหน้านี้ได้บรรยายถึงเครื่องมือ---กลไกการตรวจสอบ โครงสร้างต้นทุน และฉันทามติแบบผุดขึ้นเองของ RSN หัวข้อนี้บรรยายถึงระเบียบวิธีในการประยุกต์ใช้เครื่องมือนี้กับการเปลี่ยนผ่านทางการคลังและการปกครอง เครื่องมือนี้ใช้ได้ทั่วไป ส่วนระเบียบวิธีด้านล่างเป็นการประยุกต์ใช้ที่เป็นไปได้หนึ่งอย่าง

\subsection{โครงสร้างแรงจูงใจ}

ชื่อเสียงในฐานะทุนสร้างแรงจูงใจโดยตรงต่อพฤติกรรมที่ซื่อสัตย์ ในการดำเนินการปกติ ผู้เข้าร่วมสร้างชื่อเสียงโดยธรรมชาติผ่านธุรกรรมประจำวันและคำมั่นที่ทำได้จริง ค่าใช้จ่ายหลักคือข้ออ้าง\emph{เชิงลบ}---การบันทึกความเสียหายที่ผู้อื่นก่อขึ้น ความไม่สมมาตรนี้จูงใจพฤติกรรมที่มีประโยชน์และน่าไว้วางใจ นั่นคือ การซื่อสัตย์ไม่มีต้นทุนเพิ่มเติม ในขณะที่การก่อความเสียหายสร้างร่องรอยที่มีการจัดทำเอกสารซึ่งผู้อื่นยินดีจ่ายเพื่อรักษาไว้ ความหน้าไหว้หลังหลอก---การประกาศกฎโดยไม่ปฏิบัติตาม---เป็นรูปแบบความล้มเหลวที่แพงที่สุด เพราะแสดงให้เห็นถึงการหลอกลวงโดยเจตนา

\subsection{ระบบการคลังแบบคู่ขนาน}

องค์ประกอบสามส่วนเปิดทางให้เกิดแรงกดดันจากการแข่งขัน: (1)~\textbf{ภาษีอย่างง่าย (Simple Tax)}---อัตราตามสัดส่วนเดียวโดยไม่มีข้อยกเว้น เริ่มแรกอยู่ต่ำกว่าอัตราที่แท้จริงที่มีอยู่เพียงเล็กน้อย; (2)~\textbf{การลงทะเบียนรายจ่ายอิเล็กทรอนิกส์ (Electronic Spending Registration, ESR)}---การกลับด้านโมเดล EET ของเช็ก เพื่อให้พลเมืองสอดส่องรายจ่ายของรัฐ; (3)~\textbf{การจัดสรรภาษีที่พลเมืองกำกับ}---เริ่มต้นที่ไม่กี่เปอร์เซ็นต์ในปีแรก และไปถึง หลังจากหนึ่งทศวรรษ ที่ระดับตั้งแต่เลขสองหลักช่วงต้นไปจนถึงการโยกย้ายเต็มรูปแบบสู่ระบบที่พลเมืองกำกับ ขึ้นอยู่กับอัตราการยอมรับ จังหวะที่แท้จริงขึ้นอยู่กับความเร็วที่ทางเลือกตลาดเสรีต่อบริการของรัฐจะเกิดขึ้น และความเต็มใจของรัฐในการร่วมมือกับการเปลี่ยนผ่าน ฟังก์ชันของเวลาไม่จำเป็นต้องเป็นเชิงเส้น และไม่มีเหตุผลใดที่จะตั้งขอบเขตบนว่าการยอมรับจะไปถึงระดับใดในที่สุด

% ============================================================
% 9. MIGRATION PATH
% ============================================================
\section{เส้นทางการโยกย้าย}

การโยกย้ายดำเนินไปผ่านสามระยะ (รูปที่~\ref{fig:migration}): \textbf{ระยะที่ 1: การซ้อนทับ} (RSN ในฐานะชั้นข้อมูล รัฐเป็นผู้ให้บริการเพียงรายเดียว), \textbf{ระยะที่ 2: การแข่งขัน} (RSN มอบทางเลือกที่ใช้งานได้จริง มีการใช้งานระบบคู่) และ \textbf{ระยะที่ 3: การเปลี่ยนผ่าน} (RSN ทำได้ดีกว่าสิ่งเทียบเท่าของรัฐ การโยกย้ายโดยสมัครใจ) การเปลี่ยนผ่านย้อนกลับได้ในทุกขั้นตอน

\begin{figure}[ht]
\centering
\begin{tikzpicture}[
    phase/.style={draw, rounded corners=3pt, minimum width=1.8cm, minimum height=0.8cm, align=center, font=\scriptsize, fill=black!5},
    arr/.style={-{Stealth[length=5pt]}, thick}
]
\node[phase] (p1) at (0,0) {\textbf{ระยะที่ 1}\\การซ้อนทับ};
\node[phase] (p2) at (2.2,0) {\textbf{ระยะที่ 2}\\การแข่งขัน};
\node[phase] (p3) at (4.4,0) {\textbf{ระยะที่ 3}\\การเปลี่ยนผ่าน};

\draw[arr] (p1) -- (p2);
\draw[arr] (p2) -- (p3);
\draw[arr, dashed, gray] (p3.south) -- ++(0,-0.6) -| node[below,font=\tiny,pos=0.25] {ถอยกลับ} (p2.south);
\end{tikzpicture}
\caption{การโยกย้ายสามระยะ---ย้อนกลับได้ในทุกขั้นตอน}
\label{fig:migration}
\end{figure}

กลไกแรงกดดันหลักคือด้านการคลัง เมื่อผู้เสียภาษีแบบมีเงื่อนไขไปถึง ``ชนกลุ่มน้อยที่มีนัยสำคัญทางเศรษฐกิจ $X$''---กลุ่มที่ใหญ่พอที่การปราบปรามต่อกลุ่มนี้จะก่อความเสียหายทางเศรษฐกิจที่ไม่ใช่เรื่องเล็กน้อย และการดื้อแพ่งกลายเป็นภัยคุกคามที่น่าเชื่อถือ---รัฐก็เข้าสู่การเจรจา เพื่อเป็นตัวอย่าง เราใช้ $X \approx 15\%$ ของ GDP แต่เกณฑ์ที่แท้จริงขึ้นอยู่กับเศรษฐกิจเฉพาะแห่ง

% ============================================================
% 10. SECURITY ANALYSIS
% ============================================================
\section{การวิเคราะห์ความมั่นคง}

เราพิจารณาปฏิปักษ์ห้าประเภท ได้แก่ ผู้กระทำผิดรายบุคคล กลุ่มที่จัดตั้ง ปฏิปักษ์ระดับองค์กร ปฏิปักษ์ระดับรัฐ และปฏิปักษ์ระดับโปรโตคอล

\textbf{ความต้านทานซิบิล (Sybil resistance)}~\cite{douceur2002} การสร้างชื่อเสียงต้องอาศัยปฏิสัมพันธ์ที่ต่อเนื่องและตรวจสอบได้ ต้นทุนของการโจมตีซิบิลที่สำเร็จเพิ่มขึ้นแบบเชิงเส้นตามจำนวนอัตลักษณ์ปลอม และแบบกำลังสองตามระดับชื่อเสียงเป้าหมาย การดำเนินการหลาย DID พร้อมกันเป็นไปได้แต่แพงโดยการออกแบบ---แต่ละอัตลักษณ์ต้องสั่งสมชื่อเสียงด้วยตนเองอย่างอิสระผ่านกิจกรรมที่แท้จริง ในสังคมเสรี ต้นทุนนี้ยับยั้งการใช้ในทางที่ผิดแบบไม่จริงจัง; ในระบอบเผด็จการ DID แบบคู่ขนานกลายเป็นเครื่องมือเอาชีวิตรอด ที่เปิดทางให้เกิดการต่อต้านแบบแยกส่วน การประสานงานใต้ดิน และการเดินทางในตลาดมืดที่ปลอดภัยยิ่งขึ้น

\textbf{การป้องกันการสมรู้ร่วมคิด} รูปแบบของการรับรองซึ่งกันและกันในหมู่กลุ่มปิดสามารถตรวจจับได้ผ่านการวิเคราะห์กราฟสังคม ฟังก์ชันการรวบรวมชื่อเสียงจะลดทอนน้ำหนักข้ออ้างจากแหล่งที่รวมกลุ่มกันอย่างแน่นหนา

\textbf{ปฏิปักษ์ระดับรัฐ} การกำหนดเส้นทางแบบหัวหอม การไร้ซึ่งโครงสร้างพื้นฐานแบบรวมศูนย์ และการกระจายบทบาทผู้ตรวจสอบไปทั่วฐานผู้เข้าร่วม ทำให้มั่นใจได้ว่าการกดขี่ต้องอาศัยมาตรการที่ไม่ได้สัดส่วนกับผลประโยชน์ใดๆ

\textbf{ความเป็นส่วนตัวเทียบกับความรับผิดชอบ} ความเป็นนามแฝง---จุดกึ่งกลางระหว่างความนิรนามกับการเปิดเผยอัตลักษณ์---เปิดโอกาสให้ DID สั่งสมชื่อเสียงที่มีความหมายได้โดยไม่ต้องเปิดเผยอัตลักษณ์ในโลกจริงของผู้ถือ

% ============================================================
% 11. PRIVACY
% ============================================================
\section{ข้อพิจารณาด้านความเป็นส่วนตัว}

แบบจำลองความเป็นส่วนตัวของ RSN สร้างขึ้นบนความเป็นนามแฝง ผู้ถือควบคุมการเปิดเผยอัตลักษณ์: ขั้นต่ำ (นามแฝงล้วน) เลือกเฉพาะส่วน (เชื่อมโยงหนังสือรับรองเฉพาะ) หรือเต็มรูปแบบ (ผูกอัตลักษณ์ตามกฎหมาย) ข้ออ้างที่เผยแพร่แล้วเป็นสิ่งถาวร---ระบบรองรับการแก้ไขตามบริบทแต่ไม่รองรับการลบทิ้ง ผู้ถืออาจละทิ้ง DID ที่ถูกละเมิดและเริ่มต้นใหม่ โดยยอมสละชื่อเสียงที่สั่งสมมา

% ============================================================
% 12. RELATED WORK
% ============================================================
\section{งานที่เกี่ยวข้อง}

ข้อกำหนด W3C DID Core~\cite{did2022} และเครือข่าย Ceramic~\cite{ceramic2023} มอบรากฐานด้านอัตลักษณ์ EigenTrust~\cite{kamvar2003} ได้สาธิตการรวบรวมชื่อเสียงแบบกระจายศูนย์โดยไม่มีผู้มีอำนาจส่วนกลาง SBTs~\cite{weyl2022} ได้นำเสนอโทเคนสังคมที่โอนกันไม่ได้ RSN แตกต่างในการเน้นต้นทุนทางเศรษฐกิจในฐานะตัวกรองคุณภาพ และในกลไกการตั้งราคาความสุดโต่ง

DAOs~\cite{buterin2014} และการลงคะแนนแบบกำลังสอง~\cite{buterin2019} ได้สาธิตความเป็นไปได้ของการปกครองแบบกระจายศูนย์ RSN ได้ฉันทามติจากการเจรจาต่อรองราคาแบบสองทางแทนที่จะเป็นการลงคะแนน

ข้อเสนอนี้อาศัยทฤษฎีอนาธิปไตยทุนนิยม (anarcho-capitalist)~\cite{rothbard1973,hoppe2001} สำหรับการจัดหาบริการแบบเอกชน แต่แตกต่างในสองประเด็นสำคัญ นั่นคือ มันออกแบบเผื่อความอาฆาตและแรงกระตุ้นในการแก้แค้น แทนที่จะสมมติความมีเหตุผล และมันเสนอการเปลี่ยนผ่านที่ค่อยเป็นค่อยไปแทนการปฏิวัติ แนวคิดสัญญาประชาคมแบบผุดขึ้นเองมีที่มาจากทฤษฎีระเบียบที่เกิดขึ้นเอง (spontaneous order) ของ Hayek~\cite{hayek1973}

\textbf{อนาธิปไตยแบบอะกอริสม์ (Anarcho-agorism)} RSN มีจุดร่วมสำคัญกับเศรษฐศาสตร์ตอบโต้แบบอะกอริสม์ (agorist counter-economics)~\cite{konkin2006}---โดยเฉพาะการเน้นการสร้างสถาบันคู่ขนานและการแลกเปลี่ยนโดยสมัครใจ อย่างไรก็ตาม อะกอริสม์มีแนวโน้มที่จะสมมติว่าผลประโยชน์ส่วนตนที่มีเหตุผลเป็นกลไกการประสานงานที่เพียงพอ และมักขาดเส้นทางการโยกย้ายที่เป็นรูปธรรมจากโครงสร้างรัฐที่มีอยู่ RSN ผนวกแนวโน้มพฤติกรรมที่ไม่มีเหตุผลเข้าไว้อย่างชัดแจ้ง และมอบกลไกแรงกดดันทางการคลังสำหรับการเปลี่ยนผ่านที่ค่อยเป็นค่อยไป

\textbf{คุณธรรมนิยม (Meritocracy)} RSN อาจเป็นคำอธิบายที่ใช้งานได้จริงครั้งแรกว่าคุณธรรมนิยม~\cite{young1958} สามารถผุดขึ้นในฐานะคุณสมบัติเชิงระบบได้อย่างไร แทนที่จะเป็นเพียงคำขวัญ ระบบคุณธรรมนิยมแบบดั้งเดิมล้มเหลวเพราะต้องอาศัยผู้มีอำนาจส่วนกลางในการนิยามและบังคับใช้ ``คุณความดี'' ใน RSN คุณความดีผุดขึ้นจากการประเมินแบบสองทาง นั่นคือ ผู้ที่มอบคุณค่าจะสั่งสมชื่อเสียง; ไม่มีคณะกรรมการใดตัดสินว่าใครมีคุณความดี

\textbf{ทรัพย์สินในฐานะอภิสิทธิ์} RSN แตกต่างอย่างมีนัยสำคัญจากการปฏิบัติต่อสิทธิในทรัพย์สินของสำนักอนาธิปไตยทุนนิยมและสำนักออสเตรียนที่มองว่าเป็นสิ่งธรรมชาติและสัมบูรณ์ ในกรอบแนวคิดของ RSN ทรัพย์สินเป็น\emph{อภิสิทธิ์}---การยอมรับทางสังคมที่ในกรณีสุดขั้ว สามารถถูกเพิกถอนโดยชุมชนได้เมื่อผู้เข้าร่วมละเมิดบรรทัดฐานร่วมอย่างรุนแรง (การทรยศ การปฏิเสธที่จะปกป้องชุมชนในยามวิกฤตที่คุกคามความอยู่รอด การเอารัดเอาเปรียบอย่างเป็นระบบ) นี่ไม่ใช่การยึดทรัพย์โดยรัฐ แต่เป็นการถอนการยอมรับทางสังคมโดยเครือข่ายของความสัมพันธ์แบบสองทาง

% ============================================================
% 13. LIMITATIONS
% ============================================================
\section{อภิปรายและข้อจำกัด}

\textbf{การเริ่มต้นที่เย็นชา (Cold start)} คุณค่าของ RSN ขึ้นอยู่กับผลกระทบเครือข่าย ผู้รับเอาไปใช้ในระยะแรกเผชิญกับคุณค่าที่จำกัดจนกว่าการเข้าร่วมจะถึงเกณฑ์หนึ่ง อย่างไรก็ตาม แม้ในระยะแรกๆ เครือข่าย DID ก็ทำหน้าที่เป็นแหล่งข้อมูลเสริมที่มีประโยชน์ คาดว่าการประเมินชื่อเสียงในระยะแรกและการประเมินผู้มีอำนาจจะพัฒนาขึ้นทีละน้อยพร้อมกับความซับซ้อนที่เพิ่มขึ้น

\textbf{การต้านการเอาเปรียบและการควบคุมการเข้าถึง} DID ที่เป็นเป้าหมายอาจกำหนดค่าธรรมเนียมแบบเป็นลำดับขั้นและข้อจำกัดการเข้าถึงต่อคำสอบถามจาก DID ที่มีชื่อเสียงต่ำ นี่ไม่ใช่แบบทวิภาค แต่เป็นสเกลต่อเนื่อง นั่นคือ ยิ่ง DID ผู้สอบถามมีชื่อเสียงน้อยเท่าใด ก็ยิ่งได้รับข้อมูลน้อยลง รอนานขึ้น และจ่ายมากขึ้น กลไกนี้จูงใจการมีส่วนร่วมที่แท้จริงเหนือการบริโภคแบบเฉื่อยชา

\textbf{ความไม่เท่าเทียมในการเข้าถึง} ต้นทุนของการเผยแพร่ข้ออ้างอาจกรองข้อร้องเรียนที่ชอบธรรมออกจากผู้เข้าร่วมที่ด้อยโอกาสทางเศรษฐกิจ การบรรเทา: ผู้เข้าร่วมทุกคนทำหน้าที่เป็นผู้ตรวจสอบที่มีศักยภาพไปพร้อมกัน (หรือมอบหมายบทบาทนี้) และได้รับค่าธรรมเนียมการตรวจสอบ ในกรอบเวลาที่ยาวนานเพียงพอ ผู้เข้าร่วมที่ไม่สุดโต่งควรเข้าใกล้ความเป็นกลางทางการเงิน---รายได้จากการตรวจสอบชดเชยต้นทุนการเผยแพร่โดยประมาณ นี่คือความคาดหวังทางสถิติ มิใช่การรับประกัน

\textbf{ความไม่เท่าเทียมด้านชื่อเสียง} ผู้รับเอาไปใช้ในระยะแรกสั่งสมข้อได้เปรียบที่อาจสร้างกำแพงกั้นสำหรับผู้มาทีหลัง อย่างไรก็ตาม การประเมินชื่อเสียงดำเนินการโดยผู้ให้บริการภายนอกที่อาจเลือกบรรเทาข้อได้เปรียบของผู้มาก่อนผ่านอัลกอริทึมการรวบรวมของตน การเป็นคนแรกที่มีชื่อเสียงดีเป็นข้อได้เปรียบเชิงเปรียบเทียบทางเศรษฐกิจที่ชอบธรรม---และนั่นเป็นไปโดยการออกแบบ

\textbf{คำถามที่เปิดกว้าง} ฟังก์ชันต้นทุนที่เหมาะสมที่สุด มาตรฐานการรวบรวม การกำกับดูแลโปรโตคอล และปฏิสัมพันธ์กับข้อมูลชื่อเสียงสังเคราะห์ที่สร้างโดย AI ยังคงเป็นประเด็นสำหรับงานในอนาคต

เอกสารฉบับนี้มิได้อ้างว่า RSN จะให้ผลลัพธ์ที่เหมาะสมที่สุดในทุกสถานการณ์ แต่อ้างว่า RSN เป็นตัวแทนของทางเลือกที่\emph{เป็นไปได้}---นำไปปฏิบัติได้ในทางเทคนิค ยั่งยืนด้วยตัวเองในทางเศรษฐกิจ และเข้ากันได้ในทางสังคมกับแนวโน้มพฤติกรรมของมนุษย์ตามที่มันเป็นจริง

% ============================================================
% 14. CONCLUSION
% ============================================================
\section{บทสรุป}

เอกสารฉบับนี้ได้นำเสนอเครือข่ายสังคมชื่อเสียง ซึ่งเป็นโปรโตคอลแบบกระจายศูนย์ที่เปิดโอกาสให้เกิดการแลกเปลี่ยนชื่อเสียงแบบนามแฝงและตรวจสอบได้ หลักการความสุดโต่งราคาแพงรับประกันว่าข้ออ้างที่ฉ้อฉลจะถูกตั้งราคาจนหลุดวงโคจรไปโดยไม่ต้องเซ็นเซอร์ เส้นทางการโยกย้ายที่นำเสนอเปิดทางให้เกิดการเปลี่ยนผ่านที่ค่อยเป็นค่อยไปและย้อนกลับได้จากสถาบันที่มีอยู่เดิม การออกแบบผนวกธรรมชาติของมนุษย์ตามที่มันเป็น---รวมทั้งความใจแคบ ความอาฆาต และความปรารถนาในความสะใจจากการแก้แค้น---แทนที่จะเรียกร้องการเปลี่ยนแปลงทางอุดมการณ์ ผลลัพธ์ไม่ใช่ยูโทเปีย แต่เป็นระบบที่ความยุติธรรมเข้าถึงได้ ชื่อเสียงได้มาด้วยความพยายาม และต้นทุนของการประพฤติผิดถูกแบกรับโดยฝ่ายที่ก่อมันขึ้น

% ============================================================
% REFERENCES
% ============================================================
\bibliography{references}

\end{document}
